\documentclass[10pt,twocolumn]{article}
\usepackage{amsmath,amssymb,amsthm}
\usepackage[margin=0.85in]{geometry}
\usepackage{booktabs}
\usepackage{hyperref}
\hypersetup{colorlinks=true,linkcolor=blue,citecolor=blue,urlcolor=blue}

\theoremstyle{plain}
\newtheorem{theorem}{Theorem}
\newtheorem{proposition}[theorem]{Proposition}
\newtheorem{lemma}[theorem]{Lemma}
\newtheorem{corollary}[theorem]{Corollary}
\theoremstyle{definition}
\newtheorem{reported}[theorem]{Reported computational result}
\newtheorem{conditional}[theorem]{Conditional corollary}
\theoremstyle{remark}
\newtheorem{remark}[theorem]{Remark}

% Every displayed result that a machine check establishes names it.
\newcommand{\chk}[1]{\par\smallskip\noindent{\footnotesize\textsf{check:}
  \texttt{#1}}\par\smallskip}

\newcommand{\CF}{C_F}
\newcommand{\dg}{^{\dagger}}

\begin{document}

\title{\bfseries Nested-quotient temporal histories and homological flat bands
in strong-coupling $SU(N)$ Hamiltonian lattice gauge theory\\[2mm]
\large Master derivation, with machine-checked provenance}
\author{Alexander Smith\\ \small Independent Researcher}
\date{28 August 2026}
\maketitle

\begin{abstract}
\noindent
We study the Kogut--Susskind Hamiltonian for pure $SU(N)$ gauge theory on a
periodic cubic spatial lattice at strong coupling. After selecting the
complete gauge-invariant charge-odd one-plaquette eigenspace, fixed-order
contributions are organised as ordered magnetic histories on the Haar--Gram
quotient of reachable Wilson networks. The complete momentum dependence
through second order is the oriented edge--face Gram operator
$t_N u^2 B(k)B(k)\dg$ with $B=\partial_2\dg$ and
\[
  t_N=\frac{2N(N^2-4)}{(N^2-1)(2N^2-1)(4N^2-9)}>0,\qquad N\ge 3 .
\]
The projected spectrum is $\{0,\,t_Nu^2q,\,t_Nu^2q\}$ up to a common scalar,
with $q=4\sum_j\sin^2(k_j/2)$; the flat carrier is $\ker\partial_2$, of exact
dimension $L^3+2$.

This edition closes the one step earlier versions asserted. The channel
weights $d_\rho/N^2$ that fix $t_N$ were previously introduced by an isotropy
hypothesis; here they are \emph{derived} from the order-two Weingarten values
(Theorem~\ref{thm:weingarten}), so the all-rank second-order result is
self-contained in the strict sense. Three further statements are new: a Bloch
sum rule joining the $3\times3$ spectrum to the $(L^3+2)$-dimensional carrier
(Theorem~\ref{thm:bridge}); the identification of the vacuum-mediated route
that separates the charge-even and charge-odd hopping at every rank
(Proposition~\ref{prop:vacuum}); and an exact obstruction certificate showing
that the retained $\Gamma$/axis corpus cannot identify the disputed
fourth-order off-axis coefficient (Theorem~\ref{thm:obstruction}).

$SU(3)$ third-order and fourth-order statements remain conditional on supplied
exact-arithmetic ledgers whose raw payloads are not reproduced here, and are
labelled as such throughout. Every displayed result carries the name of the
machine check that establishes it. All statements concern a formal
finite-order projected sector: this is not a continuum glueball or mass-gap
theorem.
\end{abstract}

\section{Introduction}

Hamiltonian strong-coupling expansions organise lattice gauge dynamics around
the electric vacuum and its flux excitations. In three spatial dimensions a
fundamental plaquette excitation carries one of three face orientations, and
at second order neighbouring faces communicate through a shared link, so the
projected problem resembles a three-orbital tight-binding model. The signs of
those matrix elements are fixed by the orientation of the plaquette boundary
and by charge conjugation, and the resulting interference is the central
object here.

Flat bands generated by incidence or Gram matrices have a long history:
local topology, line-graph constructions, compact localised states, and the
need for noncontractible states on periodic systems are all established, and
closely related closed-surface structures occur in two-form spin liquids. We
do not claim that incidence flatness, cellular compact localised states, or
homological counting are new in isolation.

The content is dynamical. Starting from non-Abelian $SU(N)$ Hamiltonian
lattice gauge theory, exact strong-coupling matrix elements generate an
oriented edge--face Gram operator in the charge-odd one-plaquette sector.
Representation theory fixes its coefficient and the cubical chain complex
fixes its kernel and spectrum. The separation is exact,
\[
  \begin{aligned}
    \text{colour dynamics}&\longrightarrow t_N,\\
    \text{cellular geometry}&\longrightarrow BB\dg .
  \end{aligned}
\]

\subsection*{What is proved, and what is conditional}

Four statements are self-contained here.

\begin{enumerate}\itemsep2pt
\item For every $N\ge3$ the shared-link channel weights are $d_\rho/N^2$,
  derived from the order-two Weingarten values
  (Theorem~\ref{thm:weingarten}). This was previously an isotropy
  \emph{hypothesis}, and everything in Section~\ref{sec:second} rests on it.
\item For every $N\ge3$ the complete momentum dependence at order $u^2$ is
  $t_Nu^2BB\dg$ (Theorem~\ref{thm:allrank}).
\item The Bloch spectrum of $BB\dg$ is $\{0,q,q\}$ (Theorem~\ref{thm:fact}),
  and on the periodic $L^3$ complex the zero-mode count is $L^3+2$
  (Theorem~\ref{thm:carrier}). The two are the same statement, by a sum rule
  over the momentum grid (Theorem~\ref{thm:bridge}).
\item The retained $\Gamma$/axis corpus cannot identify the disputed
  fourth-order off-axis coefficient (Theorem~\ref{thm:obstruction}).
\end{enumerate}

Two families of statements are \emph{conditional} on supplied exact-arithmetic
derivation ledgers. Their raw class contractions, symbolic history
certificates and microscopic rank-sweep artifacts are not included here, and
repeated assertions in synthesis documents are not treated as independent
reproductions.

\begin{enumerate}\setcounter{enumi}{4}\itemsep2pt
\item For $SU(3)$, supplied ledgers report that all candidate third-order
  lifter classes vanish and report the scalar and shared-link coefficients of
  Section~\ref{sec:su3}. The census is three lifter classes evaluated over
  $32$ five-trace numerator patterns each, $96$ contractions in total, and
  nothing here checks any of them. If it is exhaustive and the linked ledger
  complete, the effective Hamiltonian stays in the incidence algebra through
  order $u^3$.
\item Supplied ledgers report the all-rank cube formula
  $\alpha_N=640/[N(N^2-1)^3]$ and its equality with the complete axial
  coefficient for $N=3,\dots,12$. Conditional on those calculations, the
  $SU(3)$ carrier acquires nonzero axial dispersion at order $u^4$. The
  complete off-axis operator remains open --- and by
  Theorem~\ref{thm:obstruction} it cannot be closed by any $\Gamma$ or axial
  datum.
\end{enumerate}

These statements concern the effective Hamiltonian projected to the degenerate
one-plaquette manifold. Virtual multi-plaquette states are retained in the
resolvents, but the external spectral problem is not the full gauge-theory
Hilbert space. At $k=0$ the three incidence branches touch and the
finite-volume separation collapses as $L^{-2}$. Nothing below proves
convergence of the strong-coupling series, persistence under unrestricted
higher-order operators, identification with the lightest physical glueball, or
a continuum mass gap.

\section{Hamiltonian and projected sector}

Let $\Lambda_L=(\mathbb{Z}/L\mathbb{Z})^3$, $L\ge3$, be a cubic spatial
lattice with periodic boundary conditions. In units of the electric coupling
the Kogut--Susskind Hamiltonian is
\begin{equation}
  H_\beta=\tfrac12\sum_\ell C_2(\ell)
  +\beta_N\sum_p\Bigl(1-\tfrac1N\operatorname{Re}\operatorname{Tr}U_p\Bigr).
\end{equation}
Dropping the additive constant,
\begin{equation}\label{eq:H}
  H=H_E-uV,\qquad V=\sum_p(\chi_p+\bar\chi_p),\qquad u=\frac{\beta_N}{2N},
\end{equation}
with $\chi_p=\operatorname{Tr}U_p$. For $SU(3)$, $u=\beta_3/6$. The
definition $u=\beta_N/(2N)$, not the $SU(3)$-specific equality, is the
all-rank convention; the archived line $Y=2\beta/3=4u$ is a definition-label
erratum and the printed coefficients are already in $u$.
\chk{the printed towers are canonical-u: 4*Delta(3u/2) reproduces them verbatim}

All quoted excitation energies use the linked, vacuum-subtracted operator
$H_{\mathrm{exc}}(u)=H(u)-E_{\mathrm{vac}}(u)I$. This changes only the scalar
ledger; incidence hopping and carrier shape are unaffected.

The electric vacuum $|0\rangle$ carries the trivial representation on every
link. For $N\ge3$ the normalised one-plaquette states of definite charge
conjugation are
\begin{equation}\label{eq:states}
  |p,\pm\rangle=\frac{\chi_p\pm\bar\chi_p}{\sqrt2}\,|0\rangle ,
\end{equation}
of unit norm by character orthogonality. Each of the four boundary links
carries the fundamental or antifundamental representation, so
\begin{equation}\label{eq:EF}
  E_{F,N}=4\Bigl(\tfrac12 \CF\Bigr)=2\CF=\frac{N^2-1}{N},
  \qquad \CF=\frac{N^2-1}{2N}.
\end{equation}
\chk{each channel gap is C\_F + C\_R/2, and the weights sum to one}
For $SU(2)$ complex conjugation is an inner gauge transformation and the
charge-odd state vanishes; this is why the construction begins at $N=3$. The
same fact appears again in Section~\ref{sec:second}: $\Lambda^2F$ is the
singlet exactly at $N=2$, which is where $t_N$ vanishes.

We work inside the charge-odd fundamental Wilson-loop cyclic sector with
trivial electric centre flux around each torus cycle. Let $P_-$ be the full
$E_{F,N}$ spectral projector there. For $L\ge3$ its range is the charge-odd
one-plaquette span, since a reduced contractible fundamental loop at that
energy has four distinct links and is an elementary plaquette. This sector
qualification also excludes the accidentally degenerate straight wrapping loop
at $L=4$. With $Q=1_--P_-$ the reduced resolvent crosses no omitted
$E_{F,N}$ eigendirection, and degenerate perturbation theory gives
\begin{equation}\label{eq:heff}
  H_{\mathrm{eff},-}=P_-H_EP_-+uH_1+u^2H_2+u^3H_3+\cdots
\end{equation}
with every $H_r$ linked, vacuum-subtracted, and carrying the appropriate
$Q$-space resolvents.

\section{Oriented incidence and the exact carrier}
\label{sec:carrier}

The periodic cubical complex has the cellular chain sequence
$C_3\xrightarrow{\partial_3}C_2\xrightarrow{\partial_2}C_1$ with
$\partial_2\partial_3=0$. One-plaquette amplitudes live in $C_2$ and their
signed leakage into links is $\partial_2c$: a cellular closed-surface
condition, not the microscopic vertex Gauss law, which the Wilson-loop basis
already satisfies through $\partial_1\partial_2=0$.

Writing $d_j=e^{ik_j}-1$ and $q(k)=\sum_j|d_j|^2=4\sum_j\sin^2(k_j/2)$, the
face-to-link Bloch incidence matrix is
\begin{equation}\label{eq:B}
  B(k)=\partial_2(k)\dg=
  \begin{pmatrix}
    d_2 & -d_1 & 0\\
    d_3 & 0 & -d_1\\
    0 & d_3 & -d_2
  \end{pmatrix}.
\end{equation}

\begin{theorem}[Incidence factorisation]\label{thm:fact}
With $w(k)=(\bar d_3,\,-\bar d_2,\,\bar d_1)^{\mathsf T}$,
\[
  BB\dg=qI-ww\dg,\qquad B\dg w=0,\qquad \|w\|^2=q,
\]
and therefore $\operatorname{spec}(BB\dg)=\{0,q,q\}$. With
$S=BB\dg-4I$ the signed shared-link adjacency,
$\operatorname{spec}S(k)=\{-4,\,q-4,\,q-4\}$.
\end{theorem}

\begin{proof}
Multiplying \eqref{eq:B} by its adjoint gives the first identity entry by
entry. Since $\|w\|^2=q$, the rank-one operator $qI-ww\dg$ annihilates $w$ and
acts as $qI$ on its orthogonal complement. Subtracting $4I$ gives the
adjacency spectrum.
\end{proof}
\chk{B B\^{}dagger = q I - d conj(d)\^{}T for the curl incidence}

For $k\neq0$ the carrier is one-dimensional and represented by $w(k)$. At
$\Gamma$, $B=0$ and all three branches meet; the normalised Bloch eigenvector
has no direction-independent limit there, which is the familiar
singular-flat-band mechanism behind the incompleteness of a purely compact
localised basis.

\begin{theorem}[Finite-torus carrier]\label{thm:carrier}
For the periodic cubic cellulation of $T^3$,
$\ker\partial_2=\operatorname{im}\partial_3\oplus H_2$ with
$\dim\operatorname{im}\partial_3=L^3-1$ and $\dim H_2=3$, so
$\dim\ker\partial_2=L^3+2$.
\end{theorem}

\begin{proof}
There are $L^3$ three-cells and no four-cells, and $b_3(T^3)=1$, so
rank$\,\partial_3=L^3-1$ by rank--nullity. Discrete Hodge decomposition gives
$\ker\partial_2=B_2\oplus H_2$ with $\dim H_2=b_2(T^3)=3$.
\end{proof}

The count is verified over $\mathbb{Z}$ rather than re-derived from the
formula: the complex is built from the boundary formulas of
Appendix~\ref{app:chain}; the exhibited cycles --- every elementary cube
boundary together with three wrapping sheets --- bound the kernel from below,
and rank over $\mathbb{F}_p$ bounds it from above, since rank can only drop
under reduction. The two bounds meet at $L=1,\dots,5$.
\chk{dim Z\_2 = L\^{}3 + 2 by rank, not by re-arranging the formula}
\chk{cube boundaries and three wrapping sheets SPAN Z\_2}

\begin{remark}
$L^3+2$ is the nullity of the down boundary Laplacian, not a Betti number; the
Betti contribution is the final $3$. The restriction $L\ge3$ belongs to the
twelve-neighbour Bloch adjacency, where $x+e$ and $x-e$ coincide at $L=2$. The
chain-level count has no such restriction and holds at $L=1,2$ as well.
\chk{the L\^{}3+2 count is chain-level, not the Bloch convention}
\end{remark}

The following sum rule joins Theorems~\ref{thm:fact} and~\ref{thm:carrier},
which are otherwise two unrelated statements: a $3\times3$ spectrum and an
$(L^3+2)$-dimensional space.

\begin{theorem}[Bloch--chain bridge]\label{thm:bridge}
On the periodic $L^3$ complex,
\[
  \sum_{k}\bigl(3-\operatorname{rank}B(k)\bigr)=L^3+2 ,
\]
the sum running over the $L^3$ allowed momenta, with nullity profile
$3$ at $\Gamma$ exactly once and $1$ at each of the other $L^3-1$ momenta.
\end{theorem}

\begin{proof}
$B(0)=0$ contributes $3$. For $k\neq0$ some $d_j\neq0$, and
Theorem~\ref{thm:fact} gives $\operatorname{rank}B(k)=2$, contributing $1$
each. Summing gives $3+(L^3-1)=L^3+2$.
\end{proof}
\chk{the Bloch and chain routes to the carrier agree}

\begin{remark}
The two decompositions produce the same two numbers from different objects.
The $3$ is the triple degeneracy of $B(0)=0$, not $b_2(T^3)$; the $L^3-1$
counts non-$\Gamma$ momenta, not cubes. Their agreement is a consistency check
on the whole construction, not a restatement.
\end{remark}

\subsection{The wrapping sheets are cycles, not harmonic}

Theorem~\ref{thm:carrier} exhibits its three $H_2$ generators explicitly as
wrapping sheets $s_{(ij)}=\sum_{x:\,x_m=0}[x;i,j]$. These satisfy
$\partial_2 s=0$ exactly, so they are cycles and they span the quotient
$\ker\partial_2/\operatorname{im}\partial_3$. They are \emph{not} harmonic:
$\partial_3\dg s\neq0$, and
\begin{equation}
  \frac{\langle s,\;\partial_3\partial_3\dg\,s\rangle}{\|s\|^2}=2
  \qquad\text{exactly, for every }L\ge2 .
\end{equation}
\chk{FINDING: the wrapping sheets are cycles but NOT harmonic}

The distinction matters because $H_2$ is used in two senses: the quotient in
Theorem~\ref{thm:carrier}, which the sheets span, and the harmonic subspace
$\ker\partial_2\cap\ker\partial_3\dg$, which they do not lie in. Nothing
downstream breaks --- Proposition~\ref{prop:rigid} rests on $B\dg w=0$, which
covers all of $\ker\partial_2$ --- but a real-space statement phrased for the
harmonic subspace does not cover the generators this construction exhibits.

\section{Exact all-rank dynamics at second order}
\label{sec:second}

Two plaquettes sharing a link leave six nonshared fundamental links after the
external contractions. On the shared link the fusion channels depend on the
relative induced orientation,
\begin{equation}\label{eq:fusion}
  F\otimes\bar F=\mathbf{1}\oplus\mathrm{Adj},
  \qquad
  F\otimes F=\Lambda^2F\oplus\mathrm{Sym}^2F ,
\end{equation}
with dimensions and quadratic Casimirs
\begin{equation}\label{eq:table}
\begin{aligned}
  d_{\mathbf 1}&=1, & C_{\mathbf 1}&=0,\\
  d_{\mathrm{Adj}}&=N^2-1, & C_{\mathrm{Adj}}&=N,\\
  d_{\Lambda^2F}&=\tfrac{N(N-1)}2, & C_{\Lambda^2F}&=\tfrac{(N+1)(N-2)}N,\\
  d_{\mathrm{Sym}^2F}&=\tfrac{N(N+1)}2, & C_{\mathrm{Sym}^2F}&=\tfrac{(N-1)(N+2)}N.
\end{aligned}
\end{equation}

The six nonshared half-links contribute $3\CF$ and the fused link $C_\rho/2$,
against an external one-plaquette energy $2\CF$, so the resolvent denominator
is $\CF+C_\rho/2$ exactly.
\chk{each channel gap is C\_F + C\_R/2, and the weights sum to one}

\subsection{The channel weights are a theorem}

Earlier editions assigned the channel projector squared norm $d_\rho/N^2$ by
an isotropy hypothesis, and every statement in this section rests on that
assignment. It is not a hypothesis.

\begin{theorem}[Shared-link channel weights]\label{thm:weingarten}
In the normalised two-plaquette state, the weight of the shared-link channel
$\rho$ is exactly $d_\rho/N^2$, for both orientation families and every
$N\ge3$.
\end{theorem}

\begin{proof}
Two steps. First the six nonshared links collapse. Each plaquette contributes
a product of three independent Haar links, and a product of independent Haar
matrices is Haar, so the amplitude is
$\operatorname{Tr}(AU)\operatorname{Tr}(BU^{\pm1})$ with $A,B$ independent
Haar and $U$ the shared link. Integrating $A$ and $B$ contributes
$\delta/N$ each and leaves a pure degree-$(2,2)$ moment of $U$.

Second, two such moments settle both families. With the order-two Weingarten
values $\mathrm{Wg}(e)=1/(N^2-1)$ and
$\mathrm{Wg}((12))=-1/(N(N^2-1))$ --- the inverse of the $S_2$ Gram matrix
$\bigl(\begin{smallmatrix}N^2&N\\N&N^2\end{smallmatrix}\bigr)$ --- the index
sums give
\begin{equation}\label{eq:moments}
  \begin{aligned}
    M_{\mathrm{direct}}&=\sum_{ijkl}\bigl\langle|U_{ij}|^2|U_{kl}|^2\bigr\rangle=N^2,\\
    M_{\mathrm{cross}}&=\sum_{ijkl}\bigl\langle U_{ij}U_{kl}\bar U_{il}\bar U_{kj}\bigr\rangle=N .
  \end{aligned}
\end{equation}
For the like family, $\operatorname{Tr}(AU)\operatorname{Tr}(BU)$ decomposes
into $\mathrm{Sym}^2\oplus\Lambda^2$ by symmetrising the two index pairs
jointly, and the two squared norms are
$(M_{\mathrm{direct}}\pm M_{\mathrm{cross}})/2$. Normalising,
\[
  n_{\mathrm{Sym}^2}=\frac{N+1}{2N}=\frac{d_{\mathrm{Sym}^2F}}{N^2},
  \qquad
  n_{\Lambda^2}=\frac{N-1}{2N}=\frac{d_{\Lambda^2F}}{N^2}.
\]
For the mixed family the singlet component of $U_{ij}\bar U_{lk}$ is
$\delta_{il}\delta_{jk}/N$, of squared norm $1$ against
$M_{\mathrm{direct}}=N^2$, so $n_{\mathbf 1}=1/N^2=d_{\mathbf 1}/N^2$ and
$n_{\mathrm{Adj}}=1-1/N^2=d_{\mathrm{Adj}}/N^2$.
\end{proof}
\chk{the shared-link weights are Weingarten, not an isotropy assumption}

\begin{remark}
The derivation imports nothing: the Weingarten pair is the inverse of the
$S_2$ Gram matrix and the index sums in \eqref{eq:moments} are explicit. It is
also independent of the rest of this program, which is why it converts the
all-rank second-order result from conditional to self-contained. Appendix~%
\ref{app:wg} records the index sums.
\end{remark}

With Theorem~\ref{thm:weingarten} the resolvent weight is
\begin{equation}\label{eq:w}
  w_\rho=-\frac{d_\rho/N^2}{\CF+C_\rho/2},
\end{equation}
and summing within the two orientation families,
\begin{align}
  A_N&=w_{\mathbf 1}+w_{\mathrm{Adj}}=-\frac{2N^3}{(N^2-1)(2N^2-1)},\label{eq:AN}\\
  B_N&=w_{\Lambda^2}+w_{\mathrm{Sym}^2}=-\frac{4N(N^2-2)}{(N^2-1)(4N^2-9)}.\label{eq:BN}
\end{align}
\chk{the four channel weights follow from dimension and Casimir}
\chk{A\_N and B\_N are the channel sums, not transcriptions}

\subsection{The shared-link law}

\begin{theorem}[All-rank shared-link law]\label{thm:allrank}
For every integer $N\ge3$,
\[
  H^{(N)}_{\mathrm{eff},-}(k,u)=E_{\mathrm{flat},N}(u)I+t_Nu^2B(k)B(k)\dg+O(u^3),
\]
with $E_{\mathrm{flat},N}$ momentum independent and
\[
  t_N=B_N-A_N=\frac{2N(N^2-4)}{(N^2-1)(2N^2-1)(4N^2-9)}>0 .
\]
The three bands through second order are $E_{\mathrm{flat},N}(u)$ and
$E_{\mathrm{flat},N}(u)+t_Nu^2q(k)$ twice.
\end{theorem}

\begin{proof}
The diagonal and within-plaquette terms are scalars in face-orbital space.
Every surviving second-order inter-plaquette process uses one shared link ---
see Proposition~\ref{prop:vacuum} for why the routes that do not are absent
from this sector --- and its relative orientation is the incidence product
appearing in $S=BB\dg-4I$. The contributions \eqref{eq:AN} and \eqref{eq:BN}
therefore assemble as $d_NI+t_NS$; absorbing $-4t_N$ into the scalar and
applying Theorem~\ref{thm:fact} gives the bands. Positivity: for $N\ge3$ every
denominator factor is positive and $N^2-4>0$.
\end{proof}
\chk{t\_N = B\_N - A\_N}\chk{t\_N > 0 for N >= 3}

Two consequences,
\begin{equation}\label{eq:tsmall}
  t_3=\frac{5}{612},\qquad
  t_N=\frac1{4N^3}-\frac1{16N^5}-\frac{77}{64N^7}-\frac{1021}{256N^9}+\cdots
\end{equation}
\chk{t\_2 = 0 and t\_3 = 5/612}\chk{large-N expansion of t\_N through 1/N\^{}9}

\subsection{The vacuum-mediated route}

The proof of Theorem~\ref{thm:allrank} asserts that only shared-link routes
survive. That is false in general and true here, for a reason worth stating:
it is a charge-conjugation selection rule, and omitting it is a recorded
error, not a hypothetical one.

\begin{proposition}[Vacuum-mediated route]\label{prop:vacuum}
The route $\langle p'|V|0\rangle\langle 0|V|p\rangle/E_0$ connects any pair of
plaquettes, sharing a link or not, and contributes $1/\CF$ per pair. It
vanishes identically in the charge-odd sector, because $V$ is charge-even and
$|0\rangle$ is charge-even, so $\langle0|V|p,-\rangle=0$. Consequently the
charge-odd hopping $t_N=B_N-A_N$ carries no such term, while the charge-even
hopping does:
\[
  \ell_N=A_N+B_N+\frac1{\CF}
  =-\frac{2N(3N^2-5)}{(N^2-1)(2N^2-1)(4N^2-9)} .
\]
\end{proposition}
\chk{ell\_N = A\_N + B\_N + 1/C\_F, the vacuum-mediated route at every rank}

\begin{remark}
At $N=3$, $A_3+B_3=-481/612$ and $1/\CF=3/4$, giving $\ell_3=-11/306$. The
first of these is exactly the value once published for the charge-even hopping
and later superseded; the correction is precisely the omitted route. Stating
the selection rule turns a one-rank erratum into an all-rank formula.
\end{remark}

\subsection{Finite-volume width}

Let $q_{\max}(L)$ be the largest sampled $q$ on the periodic $L^3$ grid. Then
\begin{equation}\label{eq:qmax}
  q_{\max}(L)=
  \begin{cases}
    12, & L\text{ even},\\[2pt]
    12\cos^2\!\bigl(\pi/2L\bigr), & L\text{ odd},
  \end{cases}
\end{equation}
because only at even $L$ does some $k_j$ reach $\pi$ exactly. The width of the
full retained charge-odd manifold is $W^{(-)}_{N,L}=q_{\max}(L)\,t_Nu^2+O(u^3)$,
so for even $L$ and in the thermodynamic limit $[u^2]W^{(-)}=12t_N\sim3/N^3$.
This concerns the stiffness of the dispersive doublet, not motion of the
lowest branch, which is flat at this order.
\chk{the zone maximum of q is 12 only at even L}

The smallest nonzero incidence eigenvalue is $q_{\min}=4\sin^2(\pi/L)$,
attained by one unit of momentum in a single direction.
\chk{q\_min on the L-torus grid is 4 sin\^{}2(pi/L)}

\subsection{The orientation signs are essential}

Removing the orientation signs --- from the entries of \eqref{eq:B} and from
the phases, so each $e^{ik}-1$ becomes $e^{ik}+1$ --- gives an unsigned
incidence $\mathcal N(k)$ with $A+4I=\mathcal N\mathcal N\dg$.
Its adjacency spectra at $\Gamma$ and $R$ are $\{12,0,0\}$ and $\{-4,-4,-4\}$,
which share no eigenvalue, so no level can be momentum independent. The flat
branch is therefore not a generic consequence of three face orbitals or cubic
symmetry.

This control is not hypothetical: the unsigned spectrum is the charge-even
sector's, and the two spans are the two bandwidths. The signed span is
$8-(-4)=12$, giving the charge-odd width $12t_-=5/51$ at $N=3$; the unsigned
span is $12-(-4)=16$, giving the charge-even width $16|\ell_3|=88/153$.
\chk{the two band spans ARE the two incidence spectra}

\section{$SU(3)$ through third order}
\label{sec:su3}

Everything in this section is conditional on supplied exact-arithmetic
ledgers. The within-plaquette charge-odd tower contributes
$8/3+u+u^2/2+7u^3/32$.

\begin{reported}[Third-order ledger]\label{rep:su3}
The supplied ledgers report, in exact rationals,
\[
  b_3=\frac{1975}{124848},\qquad
  \mathrm{leak}_3=-\frac{12331}{249696},
\]
and the assembled flat scalar
$d_3=\tfrac7{32}+12\,\mathrm{leak}_3-4b_3=-\tfrac{109151}{249696}$.
\end{reported}
\chk{d\_3 = 7/32 + 12*leak\_3 - 4*b\_3}
\chk{leak\_3 is assembled from the domino diagonal and the vacuum piece}

\begin{conditional}[Third-order factorisation]\label{cor:su3}
If the lifter census is exhaustive and the linked ledger complete, then
\[
  H_{\mathrm{eff},-}(k,u)=E_{\mathrm{flat}}(u)I+\tau(u)B(k)B(k)\dg+O(u^4),
\]
\[
  \begin{aligned}
    E_{\mathrm{flat}}(u)&=\tfrac83+u+\tfrac{11}{306}u^2-\tfrac{109151}{249696}u^3,\\
    \tau(u)&=\tfrac5{612}u^2+\tfrac{1975}{124848}u^3 .
  \end{aligned}
\]
\end{conditional}
\chk{E\_flat and t(u) carry the ledger coefficients}

The second-order assembly is $\tfrac12+12\,\mathrm{leak}_2-4t_-$ with
$\mathrm{leak}_2=-11/306$, giving the charge-odd diagonal $7/102$ and then the
flat scalar $11/306$.
\chk{d\_- = 1/2 + 12*leak\_2, and leak\_2 = -11/306}

\subsection{One assembly, both sectors, both orders}

The scalar ledger is not a list of separate facts. Writing $\lambda$ for the
adjacency eigenvalue and $\mathrm{tower}_{r,s}$ for the within-plaquette term
in canonical $u$,
\begin{equation}\label{eq:assembly}
  E_s(\lambda,r)=\mathrm{tower}_{r,s}+12\,\mathrm{leak}_{r,s}+\lambda\,t_{r,s},
\end{equation}
with $\lambda\in\{-4,8\}$ in the charge-odd sector (carrier and dispersive
top), $\lambda\in\{12,-4\}$ in the charge-even sector, and $\lambda=0$
returning the on-site diagonals. This reproduces every registered diagonal
and band value --- nine of them --- across both sectors and both orders. The
tower term is the certified coupling conversion $4\Delta(3u/2)$, so
\eqref{eq:assembly} takes no unregistered input.
\chk{one assembly formula gives every registered band value}

\begin{remark}[What the external comparison can and cannot do]
Hamer's 1989 rest-frame series agrees with $E_{\mathrm{flat}}$ at every shared
order. It cannot test $\tau$. The hopping enters the spectrum only through
$\tau(u)q(k)$ and $q(0)=0$, so every $\Gamma$-point datum is blind to it; the
one-parameter family $(b_3+\varepsilon,\ \mathrm{leak}_3+\varepsilon/3)$
leaves $d_3$ and therefore the entire comparison unchanged. What the agreement
pins is the combination $12\,\mathrm{leak}_3-4b_3$. The $\lambda=8$ band top
separates them, via \eqref{eq:assembly}.
\chk{FINDING: no Gamma-point datum can constrain the hopping}
\end{remark}

\section{What the homology protects}

\begin{proposition}[Boundary-factorised rigidity]\label{prop:rigid}
If an order-$r$ Bloch correction has the form
$H_r(k)=a_rI+b_rS(k)+B(k)M_r(k)B(k)\dg$, then for every $k\neq0$,
$H_r(k)w(k)=(a_r-4b_r)w(k)$: the correction shifts the carrier without
dispersing it.
\end{proposition}

\begin{proof}
$B\dg w=0$ and $Sw=-4w$.
\end{proof}

In real space this is the orthogonality of the two cellular Laplacians,
$L^{\downarrow}_2L^{\uparrow}_2=L^{\uparrow}_2L^{\downarrow}_2=0$, so any
polynomial in them acts on $\ker L^{\downarrow}_2\cap\ker L^{\uparrow}_2$ by
its constant term. The proposition is not an unconditional all-orders
statement: $\{BMB\dg\}$ is not a two-sided ideal in $\mathrm{End}(C_2)$, and a
general higher-order correction need not factor through links. Cubic symmetry
alone permits carrier-projected fourth-order shapes built from
\begin{equation}\label{eq:shapes}
  1,\quad q,\quad e_2,\quad \frac{4e_2}{q},\quad \frac{e_3}{q},
\end{equation}
with $e_2=\sum_{i<j}a_ia_j$, $e_3=a_1a_2a_3$ and $a_i=4\sin^2(k_i/2)$.

\section{The fourth-order boundary}
\label{sec:fourth}

Two supplied fourth-order records agree on the axial datum and on the
appearance of mobility, and disagree on one off-axis mixed-gradient
coefficient. This edition does not adjudicate them, and no fourth-order band
or bandwidth is used above. What can be settled is whether the retained corpus
\emph{could} adjudicate them.

\begin{theorem}[Obstruction certificate]\label{thm:obstruction}
The two recorded off-axis coefficients differ, in the shape basis
\eqref{eq:shapes}, by $4\Delta_C e_2$. On any one-dimensional axial cut two of
the three $a_i$ vanish, so $e_2$ is the \emph{zero polynomial} there.
Consequently no $\Gamma$-point or axial datum distinguishes the two records,
at any precision. They differ at $M$ and at $R$.
\end{theorem}
\chk{FINDING: the retained Gamma/axis data cannot identify C\_shp}

\begin{remark}
The vanishing is polynomial, not numerical: this is non-identifiability, not a
limitation of the data in hand. Re-anchoring cannot substitute, and neither
record is preferred here. An off-axis contraction is the only decider.
\end{remark}

On the axial cut itself the invariants collapse to $q=L(k)=4\sin^2(k/2)$ with
$e_2=e_3=0$, and dividing the raw cube-boundary numerator $(\alpha_N/4)L^2$ by
$\|w\|^2=q=L$ leaves a single power of $L$ with coefficient $\alpha_3/4=5/48$.
\chk{on an axial cut the mixed invariants vanish and the norm divides}

\section{Scope}

This work does not establish a four-dimensional infinite-volume Yang--Mills
mass gap, a continuum glueball mass, convergence of the strong-coupling series
to the continuum, regulator independence of the first dispersive order, the
full $SU(3)$ fourth-order rest coefficient or off-axis shape, or an
identification of any level here with a physical glueball. At $\Gamma$ the
three face orbitals transform as an axial vector of the cubic group, which
with parity even and charge conjugation odd gives the retained-space label
$T_1^{+-}$; this is an operator-space assignment, not a spectral one.

The carrier is macroscopically degenerate at finite $L$, but its nearest
incidence separation is $\Delta_L(u)=4\tau(u)\sin^2(\pi/L)+O(u^4)$, which
falls as $L^{-2}$ and so provides no volume-uniform spectral isolation.
\chk{Delta\_L = 4 tau(u) sin\^{}2(pi/L) is positive and falls as L\^{}-2}

\section{Reproducibility}

A portable standard-library Python program, \texttt{verify\_core.py}, checks
the rational coefficient ledger, the incidence identity, the chain condition
$\partial_2\partial_3=0$, the finite-torus ranks, and the Weingarten
derivation of Theorem~\ref{thm:weingarten}. It has no dependencies and runs in
under a second. Two of its checks are stronger than stated
elsewhere: the incidence identity is \emph{decided} in exact Gaussian-rational
arithmetic at rational torus points rather than checked to a residual, and the
torus ranks come with an explicit kernel basis.

Every displayed result above names the machine check that establishes it. The
checks live in the accompanying repository, where
\texttt{workhouse verify --only '<name>'} re-runs any one of them alone, with
its numbers, in about a second. Machine checks certify algebra relative to
stated definitions and inputs; they do not replace a derivation, and they do
not certify an uncompleted off-axis contraction --- which is why
Section~\ref{sec:fourth} stops where it does.

\appendix

\section{Channel weights in closed form}\label{app:closed}

Substituting \eqref{eq:table} into \eqref{eq:w},
\begingroup\small
\begin{align*}
  w_{\mathbf 1}&=-\frac{2}{N(N^2-1)}, &
  w_{\Lambda^2F}&=-\frac{N-1}{(N+1)(2N-3)},\\
  w_{\mathrm{Adj}}&=-\frac{2(N^2-1)}{N(2N^2-1)}, &
  w_{\mathrm{Sym}^2F}&=-\frac{N+1}{(N-1)(2N+3)} .
\end{align*}
\endgroup
Their family sums are \eqref{eq:AN} and \eqref{eq:BN}, and the difference has
numerator $2N(N^2-4)$, which is the $SU(2)$ zero and the positivity for
$N\ge3$ at once.

\section{The Weingarten index sums}\label{app:wg}

With $\mathrm{Wg}(e)=1/(N^2-1)$ and $\mathrm{Wg}((12))=-1/(N(N^2-1))$, and
$\sigma,\tau$ running over $S_2$,
\[
  \begin{aligned}
    &\bigl\langle U_{i_1j_1}U_{i_2j_2}
       \bar U_{i'_1j'_1}\bar U_{i'_2j'_2}\bigr\rangle\\
    &\qquad=\sum_{\sigma,\tau}\mathrm{Wg}(\sigma\tau^{-1})
      \prod_a \delta_{i_ai'_{\sigma(a)}}\delta_{j_aj'_{\tau(a)}} .
  \end{aligned}
\]
Summing the delta products over all $N^4$ index quadruples,
\[
  M_{\mathrm{direct}}=\mathrm{Wg}(e)(N^4+N^2)+2\,\mathrm{Wg}((12))N^3=N^2,
\]
\[
  M_{\mathrm{cross}}=2\,\mathrm{Wg}(e)N^3+\mathrm{Wg}((12))(N^4+N^2)=N .
\]
Both are confirmed by explicit summation at $N=3,4,5$.

\section{The $SU(3)$ coefficient ledger}\label{app:ledger}

\begin{center}
\begin{tabular}{@{}lll@{}}
\toprule
Order & Quantity & Value\\
\midrule
$u^0$ & electric plaquette & $8/3$\\
$u^1$ & scalar & $1$\\
$u^2$ & $BB\dg$ coefficient & $5/612$\\
$u^2$ & flat scalar & $11/306$\\
$u^3$ & $BB\dg$ coefficient & $1975/124848$\\
$u^3$ & per-neighbour leakage & $-12331/249696$\\
$u^3$ & flat scalar & $-109151/249696$\\
\bottomrule
\end{tabular}
\end{center}
\chk{the manuscript's SU(3) ledger is this registry, value by value}

The $u^1$ row is $SU(3)$-only rather than an $SU(3)$ specialisation: it is
$-\langle p,-|V|p,-\rangle=1$, which needs the $\epsilon$ channel and vanishes
for $N\ge4$. The other six rows specialise all-rank statements.

\section{Finite-volume chain calculation}\label{app:chain}

For $i<j$,
\[
  \partial_2[x;i,j]=[x+e_i;j]-[x;j]-[x+e_j;i]+[x;i],
\]
\begin{align*}
  \partial_3[x;1,2,3]=\;&[x+e_1;2,3]-[x;2,3]\\
  &-[x+e_2;1,3]+[x;1,3]\\
  &+[x+e_3;1,2]-[x;1,2].
\end{align*}
Substituting the second into the first cancels every edge pairwise, so
$\partial_2\partial_3=0$ over $\mathbb Z$.
\chk{d\_2 d\_3 = 0 on the built complex}
On the torus $\ker\partial_3$ is generated by the sum of all oriented cubes,
so $\operatorname{rank}\partial_3=L^3-1$.
\chk{rank d\_3 = L\^{}3 - 1 on the built complex}

\end{document}
